\chapter{Drawing a network}\label{ch:drawing}

A network in \app\ is a set of buses with components connected to them. A component's
connections are its \emph{terminals}: one for a load, shunt, generator, SVC or external
grid, two for a transformer, line or impedance. What a terminal is connected to is stored
in the component itself (its \code{bus}, \code{bus\_hv}, \code{bus\_from}\dots\ fields),
so the drawing can never disagree with the electrical model. The wires on the canvas only
show those connections.

\section{Placing components}\label{sec:placing}

\begin{minipage}[t]{0.6\linewidth}
The \gui{Library} lists every component type (\cref{ch:components}). There are two ways
to place one:

\begin{itemize}
  \item \textbf{Click and click.} Click the entry in the library, then click the canvas
    where the component should go. Hold \keys{\shift} while clicking the canvas to place
    several in a row. A right-click or \keys{Esc} drops the tool.
  \item \textbf{Drag and drop.} Drag the entry from the library onto the canvas.
\end{itemize}

Positions snap to a 10-unit grid. A new component gets an id from its type (\code{T1},
\code{LI2}, \code{LD3}, \dots) and sensible values to start with: a 10~MVA, 40/12~kV,
7\,\% transformer, a 1~km cable with $0.16+j0.08\,\Omega/\text{km}$, a 1~MW load at
cos~$\varphi=0.9$ lagging, a 2~MVAr capacitor, a 100~MW PV generator at 1~pu. New buses
are named with letters: \code{a}, \code{b}, \code{c}, \dots
\end{minipage}\hfill
\begin{minipage}[t]{0.34\linewidth}\vspace{0pt}
  \centering\shot[0.9\linewidth]{panel-library.png}
\end{minipage}

\subsection{Automatic connection}

A component connects itself when it is dropped near buses. Each terminal attaches to the
nearest bus within reach (about 95 scene units), and a bus already taken by another
terminal of the same component is skipped. So:

\begin{itemize}
  \item a transformer, line or impedance dropped \emph{between} two buses connects to both;
  \item a load, shunt, generator or SVC dropped just below a bus hangs from it;
  \item a component dropped in open space stays unconnected, and you wire it by hand.
\end{itemize}

\section{Connecting by hand}\label{sec:connecting}

Every terminal is drawn as a small circle on its symbol. To connect it:

\begin{steps}
  \item Press the mouse on the terminal and drag. While you drag, every bus the terminal
    could land on is outlined faintly, and the one under the cursor is outlined strongly.
  \item Release on a bus to connect the terminal to it. Release on empty canvas instead,
    and \app\ creates a new bus there and connects to it, in one undo step.
\end{steps}

\keys{Esc} or a right-click during the drag cancels it. To disconnect a terminal, set its
bus field to empty in \gui{Properties} or in the data manager. A component with an
unconnected terminal is allowed: it gives a warning, and the studies leave it out until
it is wired.

\section{Buses}

A bus is drawn as a bar, the \code{busbar} kind, or as a small dot, the \code{junction}
kind, set in its \gui{Kind} field. Junctions suit points inside a feeder, such as the
cable joints of S1. A bar is vertical by default.

\begin{itemize}
  \item \textbf{Orientation}: \menu{Edit > Flip bus orientation} (\keys{\ctrl + \Alt + O})
    or the bus's context menu (\gui{Switch vertical / horizontal}).
  \item \textbf{Length}: select the bus and drag the square handle at either end. The
    opposite end stays put, and the length stays a whole number of grid steps.
  \item \textbf{Where wires land}: a wire meets the bar at the point nearest to the
    component's terminal, so several branches tap one bar at different points, as in a
    real single-line diagram.
\end{itemize}

\section{Wires}

A wire leaves each terminal in the direction the symbol points, and it runs at right
angles to its bus. You can shape it with bends:

\begin{itemize}
  \item \textbf{Add a bend}: double-click the wire, or right-click it and choose
    \gui{Add bend here}.
  \item \textbf{Move a bend}: select the wire, then drag the square handle of the bend.
  \item \textbf{Remove bends}: right-click a bend handle for \gui{Remove this bend}, or
    the wire for \gui{Remove all bends}.
\end{itemize}

Bends are stored with the project and move with the drawing. Every bend edit can be
undone.

\section{Symbols and labels}

\begin{itemize}
  \item \textbf{Move}: drag a symbol, or a rubber-band selection of several. Wires follow.
  \item \textbf{Rotate}: \menu{Edit > Rotate 90 deg} (\keys{\ctrl + \Alt + R}) or the
    context menu. Labels stay upright whatever the rotation.
  \item \textbf{Labels}: every symbol shows its name and a line of key data, for example
    \emph{58.5 km}, \emph{166.5 MW 30 MVAr} or \emph{81.5 MVA 70/40 kV 8\%}. Drag a label
    to put it somewhere else. The position is saved with the project.
  \item \textbf{Out of service}: an element whose \gui{In service} box is cleared is drawn
    grey and dashed. No study takes it into account.
\end{itemize}

\section{Selecting}

Click an element to select it. Drag on empty canvas to draw a rubber band that selects
everything it touches. \keys{\ctrl} + click adds an element to the selection or removes
it, and \keys{\ctrl + A} selects all. The \gui{Properties} panel follows the selection.
Several elements of one type are edited together (\cref{sec:properties}).

\section{The property editor}\label{sec:properties}

\begin{minipage}[t]{0.55\linewidth}
\gui{Properties} shows a form for the selected element, generated from the element's
fields and grouped as in \cref{ch:components}: \gui{General}, \gui{Connection},
\gui{Parameters}, and \gui{Zero sequence} or \gui{Load flow} when the element has such
data. The internal \gui{id} is shown read-only at the top. Bindings and overrides refer
to elements by this id.

\begin{itemize}
  \item \textbf{Numbers and expressions.} A numeric field takes a number or an expression
    over the project variables, such as \code{(243 + A)/1000}. An expression is stored as a
    \emph{binding}, a link from the field to the expression kept in the project. The field
    turns amber and shows the value it evaluates to underneath, like a spreadsheet cell
    (\cref{ch:parameters}). Clear the field to remove the binding.
  \item \textbf{Validation.} A value outside a field's range (a negative length, a cos
    $\varphi$ above 1) is refused. The field turns red and its tooltip says why.
  \item \textbf{Modes.} Fields that depend on a mode appear only when they apply. Set an
    external grid's \gui{Impedance from} to \code{ik} and the \gui{Ik} and \gui{cos phi}
    fields replace \gui{Z Thevenin}.
  \item \textbf{Empty means default.} Some fields accept an empty value with a meaning of
    its own: a load's \gui{U nominal} empty means ``the bus base'', and a generator's
    \gui{Q max} empty means ``no limit''.
  \item \textbf{Several elements.} With several elements of one type selected, an edit
    applies to all of them in one undo step. A field whose values differ shows
    \emph{(several values)}. Names and bus connections are not edited in bulk.
  \item \textbf{Committing.} A field is committed when you press \keys{Enter} or leave it.
    Leaving a field without changing it does not create an undo step.
\end{itemize}
\end{minipage}\hfill
\begin{minipage}[t]{0.4\linewidth}\vspace{0pt}
  \centering\shot[0.95\linewidth]{properties-load.png}
  \captionof{figure}{A load of S1 whose \gui{P} is bound to the expression
    \code{(243 + A)/1000}.}
\end{minipage}


\section{Context menus}

\begin{minipage}[t]{0.62\linewidth}
Right-click an element for its menu. It acts on the selection the element belongs to, and
it selects the element first if it was not selected.

\begin{description}
  \item[Thevenin equivalent here] and \gui{Three-phase short circuit here} (buses only)
    run that study at this bus at once, with the study's other options at their defaults. The
    result appears as a new tab. No case is added to the project.
  \item[Properties] brings the property editor forward on this element.
  \item[Cut, Copy, Duplicate] are the clipboard commands (\cref{sec:clipboard}).
  \item[Rotate 90 deg] rotates a symbol. For a bus the entry reads \gui{Switch vertical /
    horizontal}.
  \item[In service] is checked while the element is in service. Choose it to take the
    selection out of service, or to put it back.
  \item[Delete] deletes the selection.
\end{description}
\end{minipage}\hfill
\begin{minipage}[t]{0.33\linewidth}\vspace{0pt}
  \centering\shot[0.92\linewidth]{menu-bus.png}
  \captionof{figure}{The context menu of a bus.}
\end{minipage}

Right-clicking a wire gives the bend commands described above.

\section{Clipboard}\label{sec:clipboard}

\keys{\ctrl + C} copies the selected elements, their layout and their bindings.
\keys{\ctrl + V} pastes them with fresh ids, slightly offset, and selects the pasted
elements. Connections among the copied elements are kept. A connection to a bus that was
not copied is dropped, so a pasted load whose bus stayed behind arrives unconnected.
\keys{\ctrl + X} cuts and \keys{\ctrl + D} duplicates (copy and paste in one step). The
clipboard holds plain JSON text, so it also works between two \app\ windows.

\section{Deleting}\label{sec:deleting}

\keys{Del} deletes the selection. Deleting a bus takes along the components that hang from
it. \app\ lists them and asks for confirmation first. Deleting a branch leaves its buses
in place. Bindings of deleted elements are removed with them, and undo brings everything
back.

\section{Finding an element}\label{sec:find}

\menu{Edit > Find element...} (\keys{\ctrl + F}) opens a list of every bus and component
as \emph{name (id, type)}. Type a few letters, choose an entry, or type a name or id
(either is matched, case aside) and press \keys{Enter}. The element is selected, the
diagram is centred on it, and its properties are shown. Double-clicking an id in the data
manager does the same.

\section{Undo and redo}

Every change to the project is an undoable step: adding, deleting, moving, rotating,
connecting, editing fields, wire bends, labels, the bases, the variables and the study
cases. \menu{Edit > Undo} names the step it will undo. Changes of several elements made in
one action, such as a multi-selection edit, a paste or deleting a bus with its loads, are
one step. Two things are not undoable: reloading the variable script and the annotations
toggle.

\section{Navigating the canvas}

\begin{itemize}
  \item \textbf{Zoom}: \keys{\ctrl} + mouse wheel (around the cursor), \keys{\ctrl + +} and
    \keys{\ctrl + -}. \keys{\ctrl + 0} fits the diagram, and \keys{\ctrl + 1} is actual
    size.
  \item \textbf{Pan}: drag with the middle mouse button, or hold \keys{Space} and drag
    with the left button. The mouse wheel alone scrolls.
\end{itemize}

\section{Tutorial: drawing Assignment S1 by hand}\label{sec:tutorial-s1}

This tutorial rebuilds \file{examples/s1\_a80.psn} from an empty canvas: a 40~kV grid
feeding a factory through transformer TI1 (40/12~kV), two cable sections and transformer
TI2 (12/6~kV), with loads at buses b and e.

\begin{steps}
  \item \menu{File > New}. Place five buses in a row from the library, about 240 units
    apart. They are named \code{a} to \code{e}. Set the \gui{Kind} of \code{c} and
    \code{d} to \code{junction}.
  \item Place an \gui{External grid} to the left of bus~a. It hangs from the nearest bus.
    Leave \gui{U}~=~1~pu and \gui{Impedance from}~=~\code{z\_ohm} with \gui{Z
    Thevenin}~=~0, an ideal source.
  \item Drop a \gui{Transformer} between a and b, and set \gui{S rated}~=~6~MVA, \gui{U
    rated HV}~=~40~kV, \gui{U rated LV}~=~12~kV, \gui{x}~=~7\,\%. Drop another between d
    and e: 5~MVA, 12/6~kV, 8\,\%.
  \item Drop a \gui{Line / Cable} between b and c and another between c and d. Select
    both and set, in one edit, $r=0.16\,\Omega$/km, $x=0.08\,\Omega$/km and
    $b_c=3\times10^{-6}$~S/km (type \code{3e-6}). Give the first one \gui{Parallel
    lines}~=~2.
  \item Add the variable \code{A}~=~\code{80} in the \gui{Variables} panel. Type
    \code{(217 + A)/1000} into the \gui{Length} of both cables.
  \item Place a \gui{Load} below bus~b with \gui{P}~=~\code{(117 - A)/1000}, cos
    $\varphi$~=~0.9 lagging and \gui{U nominal}~=~12~kV. Place another below bus~e with
    \gui{P}~=~\code{(243 + A)/1000} and \gui{U nominal}~=~6~kV.
  \item Open \menu{Network > Base settings...} and set the base power to 10~MVA, the base
    reference bus to a, and the base voltage there to 40~kV. Buses b, c and d get 12~kV
    and bus e gets 6~kV, by propagation through the transformer ratios.
  \item Press \keys{F5}. Without cases, every study runs with its defaults. The
    \gui{S1c) Bus voltages} tab gives $U_e = 0.9952\angle-0.5345^\circ$~pu $=
    5.9715$~kV, and the \gui{S1b)} tab gives $Z = 22.5232 + j11.1298$~pu.
\end{steps}

Save the project with \keys{\ctrl + S}. The study cases of \cref{ch:cases} then record
which buses each study looks at.
